\section*{历史注记}
\phantomsection
\addcontentsline{toc}{section}{历史注记}

（括号中的数字指本注记末尾的参考文献。）

我们已经有机会指出，平面和空间的解析几何的发展使数学家们得到 n 维空间的概念，这个概念为他们提供了一种极为方便的几何语言，可以简单而简洁地表达关于任意多个变量的方程的代数定理，特别是线性代数的所有一般结果。但是，尽管到 19 世纪中叶这种语言已为许多几何学家所惯用，它仍然纯粹是出于方便，而且由于缺乏三维以上空间的“直观”表示，在这类空间中似乎就排除了“依连续性”的论证——这类论证仅仅以“直观”为基础，在平面和（三维）空间中却是被允许的。黎曼在其关于位置分析（analysis situs）以及几何学基础的研究中，第一个通过与三维空间的类比使用了这类考虑（见第 I 章的历史注记）（*）；遵循他的榜样，许多数学家被鼓励去使用这类论证，并取得很大的成功，特别是在多个复变量的代数函数论中。但是，鉴于那个时期对空间直观的掌握非常有限，人们有理由对这类考虑的证明价值持怀疑态度，只允许在纯启发式的意义上使用它们，即认为它们只是使某些定理的真实性显得可信。因此，庞加莱在他 1887 年关于两个复变量的二重积分的残数的论文中，尽可能地避免求助于四维空间的直观：“Comme cette langue hypergéométrique répugne encore à beaucoup de bons esprits, je n'en ferai qu'un usage peu fréquent”；他为此目的所使用的技巧，使他在讨论三维空间时以拓扑论证为足，而对于这些论证，他不再迟疑求助于直观（{[}1{]}）。

此外，康托尔的发现，特别是断言 R 与 \(R^{n}\) 等势的那个著名定理{[}它似乎动摇了整个维数概念（*）{]}，表明为了把几何学和拓扑学建立在牢固的基础上，必须使它们完全摆脱对直观的一切依赖。我们已经指出（参见第 I 章的历史注记），这一需要是现代一般拓扑学观念的起源；但甚至在这一理论创立之前，就已经开始了对实 n 维空间及其直接推广（“n 维流形”）的拓扑学的严格研究，其所用的方法本属于拓扑学中称为“组合拓扑学”、或更好地称为“代数拓扑学”的那一分支。本丛书中将有一卷专门讨论拓扑学的这一分支，读者可以在那里找到关于其发展历史阶段的资料；在本章中我们只限于建立实数空间和实射影空间的最基本的拓扑性质，它们在历史上曾充当代数拓扑学方法的出发点。

